% 辅助线方法 B：过 A、D 作两条高 AE、DF
\documentclass[tikz,border=4pt]{standalone}
\usepackage{ctex}
\usepackage{amsmath}
\usetikzlibrary{calc}
\begin{document}
\begin{tikzpicture}[scale=1.0, thick]
  \coordinate (A) at (1.2,2.5);
  \coordinate (D) at (3.8,2.5);
  \coordinate (B) at (0,0);
  \coordinate (C) at (5,0);
  \coordinate (E) at (1.2,0);
  \coordinate (F) at (3.8,0);

  \draw (A)--(B)--(C)--(D)--cycle;
  \draw[red] (A)--(E);
  \draw[red] (D)--(F);
  % 直角小方块
  \draw (E) ++(0.18,0) -- ++(0,0.18) -- ++(-0.18,0);
  \draw (F) ++(-0.18,0) -- ++(0,0.18) -- ++(0.18,0);

  \node[above left]  at (A) {$A$};
  \node[above right] at (D) {$D$};
  \node[below left]  at (B) {$B$};
  \node[below right] at (C) {$C$};
  \node[below] at (E) {$E$};
  \node[below] at (F) {$F$};

  \node[below] at (2.5,-0.6) {\footnotesize 方法 B：作两条高};
\end{tikzpicture}
\end{document}
